\section{Before you start}

This manual describes every field in the settings dialog. It is organised the
way the dialog is: \textbf{General settings} first, then \textbf{Solver
settings}, then the two calibration groups, \textbf{Optimization} and
\textbf{MCMC}. Fields are referred to by the label shown on screen. Each
section ends with a table listing the fields, their defaults, and the internal
name that appears in a saved model file, in case you ever need to read one.

Most models need only the General settings. The Solver settings have defaults
that work for ordinary problems; the sections below explain when a problem is
not ordinary and which field to reach for. The Optimization and MCMC groups are
used only when you are calibrating a model against measurements, and are
ignored otherwise.

\section{General settings}

These four fields define what period is simulated and where the results go.

\paragraph{Simulation start time and Simulation end time.}
The period to simulate, in whatever time unit the model's input series use
(days, if the inflow files are in days). The start time is usually zero. The
end time deserves thought: it must be long enough for the behaviour you are
studying to actually appear. A column that has not yet broken through by the
end time does not give you a partial answer --- it gives you no answer.

\paragraph{Output File Name (All outputs).}
The file receiving the recorded results for every block and link in the model.
This is the complete record and it can become very large; see \emph{Record all
block/link outputs} under Solver settings if that becomes a problem. Leaving
this field \emph{empty} suppresses the file entirely, which is useful when you
only want the observed-data comparison.

\paragraph{Output corresponding to observed data.}
The file receiving only the quantities you have declared as Observations --- the
modelled values at the points where you have measurements. This is the file to
use when comparing the model against data, and it is far smaller than the full
output.

\section{Solver settings}

\subsection{How the time step behaves}

The solver does not use a fixed time step. It shortens the step when the
equations become hard to solve and lengthens it again when they are easy. Four
fields bound that behaviour.

\textbf{Initial time-step} is where the run begins --- but it does more than
that, and this catches people out. The step is never allowed to exceed the
initial time-step multiplied by \textbf{Maximum allowable time-step increase
factor}. That second field is a multiplier defining a ceiling, not a rate of
growth. So choosing a very small initial time-step to be safe does not simply
start the run gently: it pins the step small for the entire simulation unless
you raise the multiplier to match. Similarly \textbf{Maximum allowable
time-step decrease factor} bounds how far below the initial value the step may
fall, and \textbf{Minimum time-step} is an absolute floor.

The initial time-step has one further role: it is the spacing on which results
are written out. Making a run ten times finer therefore also produces ten times
as much output. The two cannot be set independently.

\textbf{Maximum simulation time allowed} is a real-world clock limit in
seconds. When it expires the run stops wherever it has reached. It is there so
an unattended run cannot continue indefinitely.

\subsection{How each step is solved}

Every step solves a set of nonlinear equations by repeated linear
approximation. \textbf{Newton-Raphson Tolerance} decides when the answer is
close enough.

This field is worth understanding rather than accepting. The test is made
against an absolute residual, so it interacts with the size of the numbers in
your model. In a problem where a strongly sorbing medium holds most of the mass,
the change in dissolved concentration over one step can be hundreds or thousands
of times smaller than the mass entering. If that change is smaller than the
tolerance, the solver decides it has already converged and leaves the solution
untouched --- every step, for the whole run. The run then completes and reports
success while producing a result that is entirely zero. \emph{If a model that
ought to transport something produces nothing, reduce this value first.}

When the iteration needs too many attempts the solver shortens the step by
\textbf{Newton-Raphson Time-step reduction factor}; when it fails outright it
uses the more aggressive \textbf{Newton-Raphson Time-step reduction factor when
fail}. \textbf{Maximum number of jacobian calculations allowed} is an overall
budget for the run. \textbf{Crank-Nicholson Time Weight} selects the time
weighting: $1$ is fully implicit and the most robust; $0.5$ is
Crank--Nicholson, formally more accurate but less tolerant of stiff problems.

\subsection{The Jacobian}

To take a Newton step the solver needs the sensitivity of every equation to
every unknown --- the Jacobian matrix. Two fields control how it is obtained.

\textbf{Transport Jacobian assembly} chooses how the matrix is built.
\emph{Full} re-computes every equation for each unknown it perturbs. That is
always correct but the work grows with the square of the model size.
\emph{Dependency} re-computes only the equations that the perturbed unknown can
actually influence --- its own block and its immediate neighbours --- which
produces exactly the same matrix far more quickly. Dependency is the sensible
default; use Full only if you suspect the dependency analysis itself.

\textbf{Linear solver for the Newton step} chooses what is done with the
matrix. \emph{Direct} factorises it and solves, which is the normal choice.
\emph{Inverse Jacobian} computes the inverse explicitly; that is more work and
slightly less accurate, but the inverse can be reused over several steps.
\emph{Sparse} stores only the non-zero entries, which pays off for large models
where most of the matrix is empty.

\textbf{Verify Jacobian assembly} builds the matrix both ways and reports the
largest disagreement. It roughly doubles the cost of every step and is meant for
diagnosis, not production.

\subsection{When the solution oscillates}

If the time step is coarse compared with the fastest reaction in a model, the
solution can start to oscillate: a value rises and falls on alternate steps
instead of moving steadily. Nothing is reported when this happens, because each
individual step converges perfectly well. The run finishes normally and the
results may contain excursions that cannot be real --- negative concentrations,
for instance.

\textbf{Detect and suppress solution oscillation} turns on a check for this.
It is off by default, because some models genuinely do oscillate and only you
can say whether a reversal is physical. When it is on, the solver watches the
solution and flags a variable when it reverses direction twice in a row
\emph{and} the movement is larger than \textbf{Oscillation detection tolerance}
as a fraction of that variable's own size. Both conditions matter: the reversal
pattern alone would be triggered by rounding, and the size alone by any rapidly
changing but perfectly healthy solution. The default of $0.01$ means a movement
of one percent.

\textbf{Response to detected oscillation} chooses what happens next.
\emph{Rewind} takes the solver back to the last saved state, restarts with a
step five times smaller, and throws away the results already written past that
point --- so the oscillation is removed from the output, not merely halted from
there on. \emph{Reduce} keeps the current state and simply continues with a
smaller step; it is much cheaper, but anything already written stays written.
Choose rewind when the recorded results matter most, reduce when run time does.

Either way the smaller step is held for a while, because otherwise it would
climb back to its ceiling within a few steps and the oscillation would return.
\textbf{Clean steps before relaxing the time-step ceiling} is how many
untroubled steps must pass before the solver starts allowing the step to grow
again. Oscillation is often a passing episode --- confined to a front moving
through the model --- so keeping the step small for the rest of a long run can
cost far more than the episode itself. \textbf{Maximum oscillation
interventions} limits how many times the solver will intervene before giving up
and telling you that it could not remove the oscillation, rather than refining
forever.

\subsection{Restore points}

The solver saves a complete snapshot of the model from time to time so that it
can go back to it --- after a failure, or when suppressing oscillation.

\textbf{Time steps between restore points} sets how often those snapshots are
taken. It controls two things at once: how much work is repeated when the
solver goes back (it returns to the most recent snapshot), and how much time is
spent taking snapshots. Frequent snapshots make each rewind cheap but add
overhead; infrequent ones do the opposite.

\textbf{Maximum reuses of one restore point} limits how many times the solver
may return to the same snapshot before refusing. Without such a limit a solver
could return to the same state indefinitely without ever getting past it. If
oscillation suppression reports that it gave up, raising this together with more
frequent snapshots gives it more room to work.

\subsection{What gets written, and how often}

\textbf{Record all block/link outputs} decides whether the values of every block
and link are kept. Set to \emph{No}, only the quantities you declared as
Observations are stored. This governs memory as much as file size, because
results are held in memory until the run ends: for a long run the difference can
be several gigabytes against a few hundred kilobytes. If a long run is running
out of memory, this is the field to change.

\textbf{Output write interval} is how often accumulated results are written to
disk, and \textbf{Write the outputs intermittently} decides whether that
happens during the run at all rather than only at the end. Writing during the
run lets you watch progress, at the cost of rewriting the file each time.

\textbf{Write Solution Details} produces a log of what the solver did --- where
it changed the time step, where it struggled, and why. It is the first thing to
switch on when a run behaves unexpectedly.

\textbf{Number of threads to be used} sets how much parallelism a single
simulation may use.

\section{Optimization}

These fields are used when you ask the program to find the parameter values that
best reproduce your measurements. The method is a genetic algorithm: it keeps a
collection of candidate parameter sets, scores each one by running the whole
model, and builds the next collection from the ones that scored best.

\textbf{Genetic Algorithm Population} is how many candidates exist at a time and
\textbf{Number of generations} is how many rounds are run. Multiplied together
they give the number of model runs, which is what determines how long the
calibration takes: 40 candidates over 50 generations is 2000 simulations. A
larger population searches more broadly in each round and is the better choice
when you suspect parameters trade off against one another; more generations
polish an answer that is already close.

\textbf{Cross over probability} is the chance that two selected parents are
combined rather than copied unchanged. \textbf{Mutation Probability} is the
chance that an individual value is altered at random. Mutation is what lets the
search escape a merely-good answer, so setting it to zero is unwise; setting it
far above a few percent turns the search into guesswork.

\textbf{Shaking Scale} is the size of the random adjustment applied to parameter
values, and \textbf{Shake Scale Reduction Factor} shrinks it after each
generation. Together they make the search bold at first and increasingly fine as
it settles.

\textbf{Number of threads to be used} runs several candidates at once. There is
nothing to gain from more threads than you have candidates, or than the machine
has cores.

\textbf{GA Output File Name} records every candidate of every generation with
its score --- the file to open to see whether the search actually converged.

\section{MCMC}

Optimization gives you the single best answer. MCMC gives you the range of
answers consistent with the data: it wanders through parameter space spending
time in each region in proportion to how well that region explains the
measurements. The spread of what it collects tells you how confident you can be
in each parameter, which a single best-fit value cannot.

\textbf{Number of MCMC samples} is how many steps each walker takes.
\textbf{Number of chains} is how many independent walkers there are. Several
walkers serve two purposes: they share the work, and if walkers that started in
different places end up agreeing, that is the usual evidence that the result can
be trusted. \textbf{Number of threads to be used} spreads them over cores; more
threads than chains achieves nothing.

\textbf{Number of samples to be discarded (burn out)} throws away the beginning
of each walk, which reflects where the walker happened to start rather than the
answer. \textbf{Initially purturb the samples} scatters the starting positions,
which is what makes the agreement between chains meaningful.

A step that is too large is nearly always rejected and a step that is too small
explores too slowly. The program therefore adjusts the step size automatically
to hit the \textbf{Desired acceptance rate}, the fraction of proposed moves
accepted; values between about $0.15$ and $0.3$ are usual. \textbf{Initial
perturbation factor} is the step size it begins with, and \textbf{Coefficient of
change of purtubation factor} how quickly it is adjusted.

\textbf{The interval at which the samples are saved} keeps one sample in every
so many. Consecutive samples resemble one another by the nature of the method,
so this discards very little information while keeping the file manageable.
\textbf{MCMC Output file} names it. \textbf{Continue MCMC using pre-created
file} resumes a previous run instead of starting over, which is how a run is
extended when it has not yet settled.

\textbf{Number of post-estimate realizations} is how many parameter sets are
drawn from the result and re-run, so that predictions come out as a range rather
than a single line. \textbf{Add error noise to the realizations} decides what
that range means: without it, it shows uncertainty in the parameters alone; with
it, it shows the uncertainty in an actual future measurement, which is wider and
usually the more honest thing to report. \textbf{Perform global sensitivity}
additionally reports how much each parameter contributes to the variation in the
results.
